\documentclass[11pt]{amsart}
\usepackage[T1]{fontenc}
\usepackage[utf8]{inputenc}
\usepackage{amssymb}
\usepackage[margin=1.05in]{geometry}
\usepackage[hidelinks]{hyperref}
\newcommand{\stamptime}{\the\numexpr\time/60\relax:\ifnum\numexpr\time-60*(\time/60)\relax<10 0\fi\the\numexpr\time-60*(\time/60)\relax}
\AddToHook{shipout/foreground}{\put(0.5\paperwidth,-\paperheight+0.55cm){\makebox[0pt]{\tiny\ttfamily unit/doc/\jobname.tex \ --- \ compiled \today, \stamptime}}}
\newtheorem{theorem}{Theorem}
\newtheorem{proposition}[theorem]{Proposition}
\newtheorem{lemma}[theorem]{Lemma}
\newtheorem{corollary}[theorem]{Corollary}
\theoremstyle{definition}
\newtheorem{definition}[theorem]{Definition}
\newtheorem{construction}[theorem]{Construction}
\theoremstyle{remark}
\newtheorem{remark}[theorem]{Remark}
\newcommand{\Om}{\Omega}
\newcommand{\Map}{\mathrm{Map}_*}
\newcommand{\Z}{\mathbb{Z}}
\newcommand{\code}[1]{\texttt{#1}}
\newcommand{\cT}{\mathsf{isTrunc}}
\newcommand{\cC}{\mathsf{isConn}}
\newcommand{\refl}{\mathsf{refl}}
\newcommand{\ev}{\mathsf{ev}}
\newcommand{\lp}{\mathsf{loop}}
\newcommand{\Kf}{\mathsf{Kf}}
\newcommand{\Hom}{\mathrm{Hom}}
\newcommand{\hc}[2]{[#1,#2]_*}
\newcommand{\Unl}{\mathsf{Unloop}}
\newcommand{\PMO}{\mathsf{PMO}}
\title{(H4) for the cup square: $\hc{K(\Z,2)}{K(B,4)}\cong\Hom(\Z,B)$}
\date{28 October 2026}
\begin{document}
\maketitle

\noindent\textbf{Purpose and status.} Informal proof for \S2 of
\code{doc/pi3-s2-postnikov-handoff-4.txt}:
\[\mathsf{h4}:\mathsf{KZ2Cohomology.H4}\ \mathsf{CupSquare.kappa},\]
that is, for every abelian group $B$ the map
$\Phi_B:\Hom(\Z,B)\to\hc{K_2}{K(B,4)}$, $g\mapsto[\kappa;K(g,4)]$ is an
equivalence. Written before the formal proof; the numbered statements are
meant to be declarations (\S\ref{sec:corr}). Foundations: those of the
development (univalence, images; function extensionality derived). No new
axiom; the sphere axiom is not used; no pushout, no spectral sequence. \emph{Refereed} independently (a fresh
reader given the note and the library): correct with fixes, all
incorporated here (explicit Lemma~9 section instead of \code{rt-base}; the
base point of the target of $T$; \code{conn-rt-n}; the conjugation in
Lemma~\ref{l:E-nat}; the single hypothesis $\cT\,6\,Y$; Lemma~\ref{l:j} via
pointing irrelevance).

\smallskip\noindent\textbf{The idea in one paragraph.} Write
$M_Y=\Map(K_2,Y)$. A bipointed map $c:K_2\to_*M_Y$ has a diagonal
$\Delta c=(x\mapsto(c\,x)(x)):K_2\to_*Y$, and $\kappa=\Delta c_2$ by
definition. We prove two theorems:
\begin{itemize}
\item[(A)] (\emph{every class is a diagonal}) for $Y=K(B,4)$,
  $c\mapsto[\Delta c]$ is a bijection $\Map(K_2,M_Y)\to\hc{K_2}{Y}$;
\item[(B)] (\emph{bilinear classification}) $g\mapsto c_2;K(g,4)_*$ is a
  bijection $\Hom(\Z,B)\to\Map(K_2,M_{K(B,4)})$, where $G_*$ is
  postcomposition.
\end{itemize}
Since $\Delta(c_2;K(g,4)_*)\equiv\kappa;K(g,4)$ (by definition, \S\ref{sec:H4}),
$\Phi_B$ is the composite of the two bijections, hence an equivalence.
(A) is where the information about the \emph{diagonal} is, and no factor 2
appears: the classical content is $\Delta^*(u\otimes u)=u^2$ in
$H^4(K(\Z,2))$, not the polarisation $\mu^*(u^2)=u_1^2+2u_1u_2+u_2^2$. The
proof of (A) uses DW1's Lemma~9 (the Yoneda step over the based path
space $(a:K_2)\times(a=*)$, i.e.\ the path fibration
$K(\Z,1)\to *\to K(\Z,2)$) to embed the set of classes into a set of
``multiplicativity defects'', and computes the defect of a diagonal by
one path induction (Proposition~\ref{p:star}). (B) uses only deloopings,
the circle property and the normalisation \code{cup-norm} of $c_2$.

\smallskip\noindent\textbf{The trap.} No step recovers a coefficient from
a cross term $c(x)(y)+c(y)(x)$. The invariant of a class $F$
(\S\ref{sec:A}) is computed on $F$ itself; on $F=\Delta c$ it is the
looping of $c$, so for $B=\Z/2$ the class $\kappa;K(1,4)$ is detected
(it is $\Delta(c_2;K(1,4)_*)$, and $c_2;K(1,4)_*\neq 0$ by (B)).

\section{Conventions and interfaces}\label{sec:conv}

As in \code{doc/cup-square.pdf}: \code{isTrunc n X} is ``$X$ is
$(n-2)$-truncated'', \code{isConn n X} is ``$X$ is $(n-2)$-connected'';
$K_m=K(\Z,m)$, $K^B_m=K(B,m)$ (\code{EMMaps.KK[B].K m}), pointed by $*$;
$p\cdot q$ is \code{trans p q} (diagrammatic, by $J$ on $q$, so
$p\cdot\refl\equiv p$); $\Om Y=(y_0=y_0,\refl)$; $F;G$ is \code{ptcomp F G}
($F$ first). For pointed $X,Y$: $\Map(X,Y)=(F:X\to Y)\times(F\,x_0=y_0)$,
pointed by $0=(\lambda\_.\,y_0,\refl)$.

DW1 notation (\code{DeloopMaps.Maps[A,a0,B,b0]}, here always $A=K_2$,
$a_0=*$): $\PMO=\Map(\Om A,\Om B)$, pointed by
$f_0=(\lambda q.\refl,\refl)$;
for $h:(a=a_0)\to(b=b_0)$ and $p:a=a_0$,
$D\,a\,b\,h\,p=(q\mapsto (h\,p)^{-1}\cdot h(p\cdot q),\ \ldots):\PMO$;
$\Kf\,F\,F_0\,a\,p:F\,a=b_0$ by $J$ on $p$ ($\Kf\,F\,F_0\,a_0\,\refl\equiv F_0$);
$\lp\,F=D\,a_0\,(F a_0)\,(\Kf F\,a_0)\,\refl$;
$C\,a\,b=(h:(a=a_0)\to(b=b_0))\times\prod_{p:a=a_0}(f=D\,a\,b\,h\,p)$;
$\Unl\,f=(F:\Map(A,B))\times(f=\lp\,F)$; $\mathsf{Rel}=\prod_a(b:B)\times C\,a\,b$.

Interfaces used (all in the library, or in \S2a/\S2b):
\begin{enumerate}
\item[(I1)] \code{KK}: $\cC\,(m+2)\,K^B_{m+1}$, $\cT\,(m+3)\,K^B_{m+1}$,
  pointed equivalences $e_m:\Om K^B_{m+1}\simeq_*K^B_m$; \code{KMap[A,B,g].K m}
  $=K(g,m)$ with \code{loopK}: $\lp\,K(g,m+1)=e_m;K(g,m);e_m^{-1}$;
  \code{KFun}: $K(h\circ g,m)=K(g,m);K(h,m)$.
\item[(I2)] \code{Lemma4}: \code{lemma4}, \code{cor5-map}; \code{PtSec}.
\item[(I3)] \code{DeloopMaps}: \code{lemma9} (a retract
  $\Unl\,f\lhd\mathsf{Rel}$), \code{Levels} (\code{s1, r1, s2, r2, s3, r3},
  \code{cell}, \code{cor10}), \code{cor11}, \code{Mult.cor12},
  \code{ptmap-path}, \code{E}.
\item[(I4)] \code{LoopCalc}: \code{loop-comp}
  ($\lp(F;G)=\lp F;\lp G$), \code{unique-loop}, \code{unique-loop1},
  \code{ptcomp-assoc}, \code{pt-inv-l/r}, \code{comp-const-l},
  \code{loop-const}.
\item[(I5)] \code{CircleProperty.Circle}: \code{evPt}, \code{pointed-equiv},
  \code{evPt-conj}, \code{evPt-loop}, \code{ptmap-conn}.
\item[(I6)] \code{CupSquare}: \code{PtMapLoops} ($\ev$, \code{ev-ap},
  \code{unev}, \code{ev-unev}), \code{iota}, \code{e1}, \code{e1inv},
  \code{c2}, \code{c2-loop}, \code{kappa}, \code{tM4}, \code{cup-norm}.
\item[(I7)] \code{LoopConj} (\code{conj}, \code{conj-trans},
  \code{ap-homotopy}), \code{EH.Abelian} (loops of a type pointed
  equivalent to a loop space commute), \code{PiOne.loop-mult}.
\item[(I8)] \code{H3Proof}: \code{rt-ptmap}, \code{h3}; \code{Bijection.Bij};
  \code{HomotopyClasses} (\code{HC}, \code{hcl}, \code{effective},
  \code{ind}); \code{Connected} (\code{conn-id}, \code{conn-elim}).
\end{enumerate}

\section{Three small lemmas}

\begin{lemma}[pointing irrelevance, sets]\label{l:pi-set}
If $\Map(X,Z)$ is a set and $\Om(Z,z_0)$ is connected, then
$(k,\zeta)=(k,\zeta')$ for all $k:X\to Z$, $\zeta,\zeta':k\,x_0=z_0$.
\end{lemma}
\begin{proof}
The goal is a proposition. $\zeta'=\zeta\cdot(\zeta^{-1}\cdot\zeta')$ and
$\zeta^{-1}\cdot\zeta'$ is merely $\refl$.
\end{proof}

\begin{lemma}[pointing irrelevance, circle]\label{l:pi-circle}
Let $(Z,z_0)$ be pointed with $\Om(Z,z_0)$ abelian (e.g.\ $Z$ pointed
equivalent to a loop space, \code{EH.Abelian}). Then two pointed maps
$K_1\to_*Z$ with the same underlying map are equal; the same holds for
$\Om K_2\to_*Z$ (precompose with $e_1$).
\end{lemma}
\begin{proof}
$\code{evPt}$ is injective, and $\code{evPt}(k,\zeta')=\mathsf{conj}\,\zeta'\,u
=\mathsf{conj}\,(\zeta^{-1}\zeta')\,(\mathsf{conj}\,\zeta\,u)=\mathsf{conj}\,\zeta\,u$
with $u=\mathsf{ap}_k\,g$ (\code{evPt-conj}, \code{conj-trans}; conjugation by a
loop is trivial in an abelian loop group). This is \code{CupSquare.conj-triv}
in general form.
\end{proof}

\begin{lemma}[lifting along a map with set fibres]\label{l:lift}
Let $X$ be 1-connected, $L:(M,m_0)\to_*(P,p_0)$ with
$\mathrm{fib}_L(p)$ a set for all $p$. Then $c\mapsto c;L$ is a bijection
$\Map(X,M)\to\Map(X,P)$.
\end{lemma}
\begin{proof}
For $\psi:X\to_*P$ the family $x\mapsto\mathrm{fib}_L(\psi\,x)$ is a family
of sets with a point over $x_0$ (from $L_0$ and $\psi_0$). Its pointed
sections are contractible (\code{lemma4}, $c=2$, $t=0$). A pointed
section $\sigma$ gives $c=(\lambda x.(\sigma x)_1,\mathsf{ap}_{\mathrm{pr}_1}\sigma_0)$
with $c;L=\psi$ (homotopy $(\sigma x)_2$, coherence from $\sigma_0$); two
lifts $c,c'$ of the same $\psi$ give two pointed sections, equal by
contractibility, so $c=c'$.
\end{proof}

\section{Classes as deloopings of zero}\label{sec:P}

Fix $Y=K^B_4$, $A=K_2$; $\PMO_Y=\Map(\Om K_2,\Om K^B_4)$, $f_0$ its base point.

\begin{lemma}\label{l:pmo}
$\PMO_Y$ is 1-connected ($\cC\,3$) and 2-truncated ($\cT\,4$); each
$\Unl\,f$ ($f:\PMO_Y$) is a set; $M_Y$ is 2-truncated.
\end{lemma}
\begin{proof}
$\PMO_Y$ is a retract of $\Map(K_1,K^B_3)$ (\code{H3Proof.rt-ptmap} with
$e_1$, $e_3$), which is equivalent to $\Om K^B_3\simeq K^B_2$
(\code{pointed-equiv}, $e_2$), 1-connected; connectedness passes to
retracts at every level (\code{conn-rt-n}, new: the library's
\code{conn-rt} is for $\cC\,2$ only). Truncation: \code{Levels.tPMO} ($c=1$, $t=2$). $\Unl\,f$ is a set
by \code{Levels[1,2].cor10}: $K_2$ is $\cC\,3$ and $K^B_4$ is
$\cT\,6=\cT(2+2+2)$. $M_Y$: \code{cor5-map} as in \code{CupSquare.tM4}.
\end{proof}

\begin{proposition}\label{p:P}
$\pi:\Unl\,f_0\to\hc{K_2}{Y}$, $(F,w)\mapsto[F]$, is a bijection.
\end{proposition}
\begin{proof}
Surjective: for a class $[F]$ (\code{HC} induction into propositions),
$\PMO_Y$ is connected, so merely $w:f_0=\lp F$. Injective: if $[F]=[F']$
then merely $e:F=F'$ (\code{effective}); the goal $(F,w)=(F',w')$ is a
proposition ($\Unl\,f_0$ is a set), and $(F,w)=(F',w\cdot\mathsf{ap}_\lp e)$;
two elements of $f_0=\lp F'$ are merely equal because this type is
connected ($\PMO_Y$ is 1-connected, \code{conn-id}). Conclude with
\code{Bijection.Bij}.
\end{proof}

\section{Theorem A: every class is a diagonal}\label{sec:A}

Here $X=K_2$ and $Y$ is any pointed type with $\cT\,6\,Y$. Everything
used follows: $\PMO$ and $M$ are $\cT\,4$; each $\Unl\,f$ and each
$(b:Y)\times C\,a\,b$ is a set (\code{Levels[1,2]}: \code{cor10},
\code{cell}); \code{cor11} applies to $K_2\to_*\PMO$; $\Om(\PMO,f_0)$ is a
loop space. Proposition~\ref{p:star} and Lemma~\ref{l:equiv} hold for any
connected $A$ in place of $K_2$. Write $L=\lp:M\to\PMO$ and
$\lambda_0:L\,0=f_0$ (\code{loop-const}); $\hat L=(L,\lambda_0):M\to_*\PMO$.

\begin{definition}[diagonal]
For $c:K_2\to_*M$ with pointing $c_0:c\,a_0=0$:
$\Delta c=(\lambda x.\,(c\,x)_1\,x,\ (c\,a_0)_2):K_2\to_*Y$. Thus
$\kappa\equiv\Delta c_2$.
\end{definition}

\subsection*{The invariant}
We write out the section of Lemma~9 (its library form goes through
\code{rt-base}, i.e.\ \code{Torsor.eqind}, which has no computation rule,
so it is not used):
\[s(F,w)\,a=\bigl(F a,\ \Kf F\,a,\ \lambda p.\,\mathsf{cell}_{F,w}(a,p)\bigr),\qquad
\mathsf{cell}_{F,w}(a,p)=J(w):f_0=D\,a\,(Fa)\,(\Kf F\,a)\,p,\]
$\mathsf{cell}_{F,w}(a_0,\refl)\equiv w$, with the retraction
$\mathsf{ret}\,\rho=((\lambda a.(\rho a)_1,\ H\,a_0\,\refl),\ (\rho a_0)_{2,2}\refl\cdot\mathsf{ap}_D(k)^{-1})$,
$H=\lambda a.(\rho a)_{2,1}$, $k$ the $a_0$-component of \code{kf-ev}~$H$;
$\mathsf{ret}(s\,u)=u$ because \code{kf-ev} at a $J$-generated $H$ is $\refl$
($J$, then $\mathsf{funext}(\lambda\_.\refl)=\refl$). So $s$ is injective. Evaluation at $a_0$ is injective on
$\mathsf{Rel}$: each $(b:Y)\times C\,a\,b$ is a set (\code{Levels.cell},
$t=2$) and $K_2$ is connected (\code{conn-elim} for the propositions
$\rho\,a=\rho'\,a$). The proof of \code{Levels.cell0} transforms
$(b:Y)\times C\,a_0\,b$ injectively into pointed sections:
$T=\mathsf{tr}_2\circ s_3\circ s_2\circ\mathsf{tr}_1\circ s_1$, where
$s_1,s_2,s_3$ have retractions and $\mathsf{tr}_1,\mathsf{tr}_2$ are the
transports to the centres of the singletons $\mathsf{Sing}=(b:Y)\times(b=b_0)$
and $\mathsf{Sing2}=(h:\PMO)\times(f_0=h)$; the target is
$\mathsf{PtSec}=\mathsf{PtSec}(\Om K_2,\refl,p\mapsto (f_0=D\,a_0\,b_0\,f_{0,1}\,p),\ldots)$.
Let $T'$ be $T$ followed by forgetting the pointing; $T'$ is injective by
Lemma~\ref{l:pi-circle}: since $D\,a_0\,b_0\,f_{0,1}\,p$ is $f_0$ by
definition, the target is $\Map(\Om K_2,(f_0=f_0,\ast))$ with
$\ast=\refl\cdot E\,f_0$ (not $\refl$); loops at $\ast$ commute by
\code{EH.Abelian} applied to the re-pointing equivalence
$(\Om\PMO,\refl)\simeq_*(f_0=f_0,\ast)$, $q\mapsto q\cdot\ast$.
Explicitly $T'(b,h,k)(p)=\bar k(p)\cdot\mathsf{ap}_{D(-,p)}(\omega^{-1})$
where $\bar k$ is $k$ after $\mathsf{tr}_1$ and
$\omega=\bar k(\refl)\cdot E^{-1}$ (transport in a family of path types,
$J$).
\[\Theta=T'\circ\mathsf{ev}_{a_0}\circ s:\ \Unl\,f_0\to\textstyle\prod_{p:\Om K_2}(f_0=f_0)\quad\text{is injective.}\]

\subsection*{The action of defects}
For $\varphi:\prod_{p:\Om K_2}(f_0=f_0)$ let
$\varphi\rhd(b,h,k)=(b,h,\lambda p.\,\varphi(p)\cdot k(p))$ on
$(b:Y)\times C\,a_0\,b$, and $(\varphi\cdot\theta)(p)=\varphi(p)\cdot\theta(p)$.

\begin{lemma}[equivariance]\label{l:equiv}
If $\varphi(\refl)=\refl$ then $T'(\varphi\rhd v)=\varphi\cdot T'(v)$.
\end{lemma}
\begin{proof}
$s_1,s_2$ commute with the action by definition; $\mathsf{tr}_1$ acts on
the $k$-component by postcomposition, which commutes with precomposition
by $\varphi(p)$ ($J$ on the path in $\mathsf{Sing}$). In the formula for
$T'$, $\omega$ becomes $\varphi(\refl)\cdot\omega=\omega$.
\end{proof}

\subsection*{The defect of a diagonal}

\begin{lemma}[$\Kf$ of a diagonal]\label{l:K}
For $c:K_2\to_*M$, $F=\Delta c$, all $a$ and $p:a=a_0$:
$\Kf\,F\,F_0\,a\,p=\Kf\,(c a)_1\,(c a)_2\,a\,p$.
\end{lemma}
\begin{proof}
$J$ on $p$ (the map $c\,a$ varies with $a$ in the motive); at
$(a_0,\refl)$ both sides are $(c\,a_0)_2$ by definition.
\end{proof}
Let $K_a:\Kf F\,a=\Kf (c a)\,a$ be its function extensionality. For any
$m:M$ let $\mu_m(a,p):\lp\,m=D\,a\,(m\,a)\,(\Kf m\,a)\,p$ be $J(\refl)$
(the canonical multiplicativity of $\lp\,m$); then
$\mathsf{cell}_{m,w}(a,p)=w\cdot\mu_m(a,p)$ ($J$; at $\refl$ both are $w$).

\begin{construction}\label{c:R}
For $c$, $F=\Delta c$, and $p:a=a_0$,
\[R(a,p)=\bigl(\mathsf{ap}_L(\Kf c\,a\,p)\cdot\lambda_0\bigr)^{-1}\cdot\mu_{c a}(a,p)
\cdot\bigl(\mathsf{ap}_{D\,a\,(F a)\,(-)\,p}\,K_a\bigr)^{-1}\ :\ f_0=D\,a\,(Fa)\,(\Kf F\,a)\,p,\]
and $w_c=R(a_0,\refl)$; so $\Delta' c=(\Delta c,w_c):\Unl\,f_0$.
\end{construction}

\begin{proposition}[the defect of a diagonal]\label{p:star}
For every $c:K_2\to_*M$,
$\ \mathsf{ev}_{a_0}(s(\Delta'c))=\varphi_c\rhd v_{00}$, where
$v_{00}=\mathsf{ev}_{a_0}(s(0,(\refl\cdot\lambda_0)^{-1}))$ does not depend on $c$ and
\[\varphi_c(p)=\bigl(\mathsf{ap}_L(\Kf c\,a_0\,p)\cdot\lambda_0\bigr)^{-1}\cdot\bigl(\mathsf{ap}_L(c_0)\cdot\lambda_0\bigr).\]
Hence $\Theta(\Delta'c)=\varphi_c\cdot\Theta(0,(\refl\cdot\lambda_0)^{-1})$
(Lemma~\ref{l:equiv}; $\varphi_c(\refl)=\refl$ by \code{trans-sym-l}).
\end{proposition}
\begin{proof}
(i) $\mathsf{cell}_{\Delta c,w_c}(a,p)=R(a,p)$ for all $(a,p)$, by $J$: at
$(a_0,\refl)$ both are $w_c$ by definition. (ii) So
$\mathsf{ev}_{a_0}s(\Delta' c)=(b,\Kf F\,a_0,\lambda p.R(a_0,p))$ with
$b=(c\,a_0)_1a_0$. Transport along $K_{a_0}$ in the $h$-component
cancels the last factor of $R$: this equals
$(b,\Kf m\,a_0,\lambda p.\,(\mathsf{ap}_L(\Kf c\,a_0\,p)\lambda_0)^{-1}\mu_m(a_0,p))$,
$m=c\,a_0$. (iii) Factor out $\varphi_c(p)$: the rest is
$(b,\Kf m\,a_0,\lambda p.\,w_m\cdot\mu_m(a_0,p))=\mathsf{ev}_{a_0}s(m,w_m)$
with $w_m=(\mathsf{ap}_L(c_0)\lambda_0)^{-1}$. (iv) $J$ on
$c_0:m=0$ (with $m$ free): $\mathsf{ev}_{a_0}s(m,(\mathsf{ap}_L c_0\cdot\lambda_0)^{-1})=v_{00}$.
\end{proof}

\begin{lemma}\label{l:phi}
$\varphi_c(p)=\bigl((\lp(c;\hat L))_1\,p\bigr)^{-1}$; so $\varphi_c=\varphi_{c'}$ iff
$\lp(c;\hat L)=\lp(c';\hat L)$ (as pointed maps, Lemma~\ref{l:pi-circle}).
\end{lemma}
\begin{proof}
\code{kf-ap}: $\Kf c\,a_0\,p=\mathsf{ap}_c\,p\cdot c_0$; $\mathsf{ap}_L$ of a
composite; \code{loop-ap}: $\lp(c;\hat L)\,p=\mathsf{conj}\,(\mathsf{ap}_Lc_0\cdot\lambda_0)\,(\mathsf{ap}_{L\circ c}\,p)$.
\end{proof}

\begin{lemma}\label{l:bij-loop}
$c\mapsto\lp(c;\hat L)$ is a bijection $\Map(K_2,M)\to\Map(\Om K_2,\Om(\PMO,f_0))$.
\end{lemma}
\begin{proof}
$\lp:\Map(K_2,\PMO)\to\Map(\Om K_2,\Om\PMO)$ has contractible fibres
(\code{cor11}, $c=1$, $\PMO$ is $\cT\,4$), and $c\mapsto c;\hat L$ is a
bijection by Lemma~\ref{l:lift} (fibres of $L$ are $\Unl$'s, sets).
\end{proof}

\begin{theorem}[A]\label{t:A}
$\Delta':\Map(K_2,M)\to\Unl\,f_0$ is a bijection; with
Proposition~\ref{p:P}, $c\mapsto[\Delta c]$ is a bijection
$\Map(K_2,M_{K^B_4})\to\hc{K_2}{K^B_4}$.
\end{theorem}
\begin{proof}
Let $z=\Theta(0,(\refl\lambda_0)^{-1})$. Injective: $\Delta'c=\Delta'c'$ gives
$\varphi_c\cdot z=\varphi_{c'}\cdot z$, so $\varphi_c=\varphi_{c'}$
(cancel pointwise), so $c=c'$ (Lemmas~\ref{l:phi}, \ref{l:bij-loop}).
Surjective: for $u:\Unl\,f_0$ let $\psi(p)=\Theta(u)(p)\cdot z(p)^{-1}$; it
is pointed ($\Theta(u)(\refl)$ and $z(\refl)$ are both the base point of
$\mathsf{PtSec}$ up to the pointings); by Lemma~\ref{l:bij-loop} there is
$c$ with $\lp(c;\hat L)=\psi^{-1}$ (pointwise inverse), so
$\varphi_c=\psi$ and $\Theta(\Delta'c)=\Theta(u)$, hence $\Delta'c=u$.
\end{proof}

\section{Theorem B: the bilinear classification}\label{sec:B}

For a pointed $G:Y\to_*Y'$ let $G_*:M_Y\to_*M_{Y'}$, $m\mapsto m;G$,
pointed by \code{comp-const-l}. For $c:K_2\to_*M_Y$ let
\[\mathcal E(c)=\ev_{\Om Y}\bigl(\code{evPt}(e_1^{-1};\lp\,c;\ev_Y)\bigr):K_2\to_*\Om^2Y,\]
the composite of looping, the mapping-loop equivalences and the circle
property, bracketed as \code{c1bar}, $e_1^{-1};(\lp c;\ev_Y)$, so that $\mathcal E(c_2)$
is by definition the left side of \code{cup-norm}: $\mathcal E(c_2)=\iota$.

\begin{lemma}\label{l:E-inj}
$\mathcal E$ is injective (for $Y=K^B_4$).
\end{lemma}
\begin{proof}
$\lp$ is injective on $\Map(K_2,M_Y)$ (\code{unique-loop}, $M_Y$ is
$\cT\,4$); pre- and postcomposition with pointed equivalences, \code{evPt}
and $\ev$ are injective.
\end{proof}

\begin{lemma}[naturality]\label{l:E-nat}
$\mathcal E(c;G_*)=\mathcal E(c);\lp^2G$ for pointed $G:Y\to_*Y'$, when
$\Map(K_2,\Om^2Y')$ is a set and $\Om^3Y'$ is connected (e.g.\ $Y'=K^{B}_4$).
\end{lemma}
\begin{proof}
By Lemma~\ref{l:pi-set} it suffices to compare underlying maps.
$\lp(c;G_*)=\lp c;\lp G_*$ (\code{loop-comp}); \code{evPt} is natural:
$\code{evPt}(\psi;N)=\lp N(\code{evPt}\,\psi)$ (\code{evPt-loop},
\code{loop-comp}). With $N=\ev_Y^{-1};\lp G_*;\ev_{Y'}$ (\code{ev-unev})
and $\ell=\code{evPt}(e_1^{-1};\lp c;\ev_Y)$ we must show, for each $x$,
$\mathsf{ap}_{\mathrm{ev}_x}(\lp N\,\ell)=\lp^2G\,(\mathsf{ap}_{\mathrm{ev}_x}\ell)$
(\code{ev-ap} on both sides). As in \code{cup-norm}: take
$\mathrm{ev}_x:\Map(K_2,\Om Y')\to_*\Om Y'$ (pointed by $\refl$), so both sides
are loopings, $\lp(N;\mathrm{ev}_x)\,\ell$ and $\lp(\mathrm{ev}_x;\lp G)\,\ell$
(\code{loop-comp}). The pointed maps $N;\mathrm{ev}_x$ and
$\mathrm{ev}_x;\lp G$ are homotopic (unpointed), by \code{ev-ap},
\code{loop-ap}, an ``$\mathsf{ap}$ of $\mathsf{conj}$'' lemma ($J$),
$\mathsf{ap}_{\mathrm{ev}_x}\pi_G=\refl\cdot G_0$ (\code{happly-funext}) and
$\ev^{-1}\circ\ev=\mathrm{id}$ (from injectivity of $\ev$; \code{unev} is
locked). By \code{LoopConj.LoopHomotopy} their loopings differ by
conjugation by a loop of $\Om Y'$ at $\refl$, i.e.\ a 2-loop, acting on a
2-loop: trivial by \code{EH.eh}.
\end{proof}

\begin{lemma}[H2]\label{l:H2}
For abelian $A,B$, $g\mapsto K(g,1)$ and $g\mapsto K(g,2)$ are bijections
$\Hom(A,B)\to\Map(K^A_1,K^B_1)$, $\to\Map(K^A_2,K^B_2)$.
\end{lemma}
\begin{proof}
Level 1. $\lp K(g,1)=\mathsf{conj}_0(g)=\mathsf{unband}\circ g\circ\mathsf{band}$
(\code{loopK 0}, \code{c0-conj}). Injective: equal maps have equal loopings,
so $g\circ\mathsf{band}=g'\circ\mathsf{band}$, and $\mathsf{band}$ is surjective.
Surjective: for $F$, $g_F=\mathsf{band}\circ(\lp F)_1\circ\mathsf{unband}$ is a
homomorphism (\code{loop-mult}, \code{band-hom}); $\lp K(g_F,1)=\lp F$
pointwise (\code{band-unband}, \code{unband-band}), pointings in a set;
so $K(g_F,1)=F$ (\code{unique-loop1}, $c=0$).
Level 2. $\lp K(g,2)=e_1;K(g,1);e_1^{-1}$ (\code{loopK 1}); conjugation by
$e_1$ is a bijection $\Map(K^A_1,K^B_1)\to\PMO$; $\lp$ is injective
(\code{unique-loop}, $c=1$) and every $\PMO$ element is a looping
(\code{cor11}).
\end{proof}

\begin{lemma}[$\Om^2$ of $K(g,4)$]\label{l:j}
Let $j_B=\lp e_3;e_2:\Om^2K^B_4\to_*K^B_2$ (so $\iota=j_\Z^{-1}$). Then
$\iota;\lp^2K(g,4)=K(g,2);j_B^{-1}$.
\end{lemma}
\begin{proof}
$\Map(K_2,\Om^2K^B_4)$ is a set and $\Om^3K^B_4$ is connected, so by
Lemma~\ref{l:pi-set} it suffices to compare underlying maps after
postcomposing with $j_B$: \code{loopK} at 3 and 2, \code{loop-comp},
\code{pt-inv-r} and \code{Puppe.loop-id} pointwise. ($j_\Z$ is by definition
\code{iota-eq.1}.)
\end{proof}

\begin{theorem}[B]\label{t:B}
$g\mapsto c_2;K(g,4)_*$ is a bijection $\Hom(\Z,B)\to\Map(K_2,M_{K^B_4})$.
\end{theorem}
\begin{proof}
$\mathcal E(c_2;K(g,4)_*)=\iota;\lp^2K(g,4)=K(g,2);j_B^{-1}$
(Lemmas~\ref{l:E-nat}, \ref{l:j}, \code{cup-norm}). Injective: by
Lemma~\ref{l:H2} and injectivity of $-;j_B^{-1}$. Surjective: for $c$,
$\mathcal E(c);j_B=K(g,2)$ for some $g$ (Lemma~\ref{l:H2}), so
$\mathcal E(c)=\mathcal E(c_2;K(g,4)_*)$ and $c=c_2;K(g,4)_*$
(Lemma~\ref{l:E-inj}).
\end{proof}

\section{(H4), its inverse, and the coefficients}\label{sec:H4}

\begin{theorem}\label{t:H4}
$\mathsf{h4}:\mathsf{H4}\ \kappa$.
\end{theorem}
\begin{proof}
By definition $\Delta(c_2;G_*)=(\lambda x.\,G((c_2x)x),\ \mathsf{ap}_G(c_2a_0)_2\cdot G_0)
\equiv\kappa;G$. Hence $\Phi_B=[-]\circ\Delta\circ(g\mapsto c_2;K(g,4)_*)$,
a composite of bijections of sets (Theorems~\ref{t:B}, \ref{t:A}); \code{Bijection.Bij}
turns it into \code{isEquiv}, i.e.\ an inverse $\Psi_B$ with both equations.
\end{proof}

\noindent\textbf{(a) The inverse.} $\Psi_B$ is, explicitly, the composite
of the inverses: a class $[F]$ is represented by $(F,w)$ (Prop.~\ref{p:P});
$c$ is the bipointed map with $\lp(c;\hat L)=(\Theta(F,w)\cdot z^{-1})^{-1}$
(Lemma~\ref{l:bij-loop}: delooping by \code{cor11}, then lifting); and $g$
is $K(-,2)^{-1}(\mathcal E(c);j_B)$, that is, with $G=e_1^{-1};\lp(\mathcal E(c);j_B);e_1:K_1\to_*K^B_1$,
$g=\mathsf{band}\circ(\lp G)_1\circ\mathsf{unband}$ (Lemma~\ref{l:H2}).
Both equations $\Psi_B\Phi_B=\mathrm{id}$, $\Phi_B\Psi_B=\mathrm{id}$ are the
\code{Bij} equations.

\noindent\textbf{(b) The generator is the constructed $\kappa$.} It
enters exactly twice: $\kappa\equiv\Delta c_2$ (definition), and
$\mathcal E(c_2)=\iota$, which is \code{cup-norm} (itself from
\code{iota}, \code{c1-ev}, \code{c2-loop}, \code{ev-ap} through
$\mathcal E$'s definition).

\noindent\textbf{(c) Coefficients.} $B$ is arbitrary (including $\Z/2$);
naturality: for $h:B\to B'$, $[\Phi_B(g);K(h,4)]=\Phi_{B'}(h\circ g)$ by
\code{KFun} ($K(h\circ g,4)=K(g,4);K(h,4)$) and associativity of
composition of classes. $\Phi_B$ is the map of \code{KZ2Cohomology.Phi}
(it uses \code{KMap[Zab,B,g].K four}).

\noindent\textbf{Remark (the trap).} The invariant $\Theta$ is not
quadratic in disguise: for $B=\Z/2$ and $g=1$, Theorem~B gives
$c_2;K(1,4)_*\neq0$, and Theorem~A gives $[\kappa;K(1,4)]\neq0$; a cross-term
invariant would see $2\cdot1=0$.

\section{Dependency audit (requirement d)}\label{sec:audit}

\noindent\textbf{No Gysin or transgression sequence is used.} The role of
the path fibration $K(\Z,1)\to *\to K(\Z,2)$ is played by DW1's Lemma~9:
the cell of a delooping is defined over the based path space
$(a:K_2)\times(a=a_0)$ by $J$, and Proposition~\ref{p:star}(i) is a $J$ over
it. No pushout, no suspension, no join, no truncation beyond
\code{HC} (set truncation, from \code{Quotient}).

\noindent\textbf{Levels.} $K_2$: $\cC\,3$ (1-connected); $K^B_4$: $\cC\,5$,
$\cT\,6$. $\Unl\,f$: \code{Levels[c=1,t=2]}, sets ($\cT(2+2+2)$).
$\PMO$: $\cT\,4$ (\code{tPMO}), $\cC\,3$ (Lemma~\ref{l:pmo}). $M$: $\cT\,4$
(\code{cor5-map 4 2}). \code{cor11} at $c=1$ for $K_2\to_*M$ and
$K_2\to_*\PMO$ ($\cT\,4$); \code{unique-loop} at $c=1$; \code{unique-loop1}
at $c=0$ ($K^B_1$ is $\cT\,3$); \code{lemma4} with $c=2,t=0$ (sets over a
1-connected base) in Lemma~\ref{l:lift}. $\Map(K_2,\Om^2K^B_4)$ is a set
(\code{cor5-map 2 2}), $\Om^3K^B_4$ is connected (\code{conn-id} three
times from $\cC\,5$).

\noindent\textbf{Library lemmas used:} (I1)--(I8). \textbf{New lemmas:} \code{conn-rt-n}; the
re-pointing pointed equivalence for \code{EH.Abelian}; the explicit
section/retraction of Lemma~9; transport in families of path types (the
action; $\mathsf{tr}_1,\mathsf{tr}_2$); lifts as pointed sections
(Lemma~\ref{l:lift}); ``$\mathsf{ap}$ of $\mathsf{conj}$'';
$\ev^{-1}\circ\ev=\mathrm{id}$; postcomposition on classes for (c);
Lemmas~\ref{l:pi-set}, \ref{l:pi-circle}, \ref{l:lift}, \ref{l:pmo},
Prop.~\ref{p:P}; injectivity of $\Theta$; Lemma~\ref{l:equiv},
\ref{l:K}, the $\mu$/cell lemma, Construction~\ref{c:R},
Prop.~\ref{p:star}, Lemmas~\ref{l:phi}, \ref{l:bij-loop}, Theorem~A;
Lemmas~\ref{l:E-inj}, \ref{l:E-nat}, \ref{l:H2}, \ref{l:j}, Theorem~B;
Theorem~\ref{t:H4}. No library file is modified.

\section{Correspondence with the formal proof}\label{sec:corr}
\emph{Status: checked} (68 modules, 0 failed; 2771 declarations, 1.1\,s;
\code{h4.opens} generated from \code{--opened}, not yet tight).
\begin{itemize}
\item \code{H4Classes}: \code{conn-rt-n}, Lemma~\ref{l:pi-set} \code{pi-set};
  \code{ClassesAt[B]}: Lemma~\ref{l:pmo} (\code{pmo-conn}, \code{unloop-set}),
  Prop.~\ref{p:P} (\code{cls}, \code{cls-inj}, \code{cls-surj}, \code{cls-equiv}).
\item \code{H4Diagonal}: \code{Section} (the Lemma~9 section \code{sec},
  retraction \code{ret}, \code{sec-inj}); \code{RelEval.ev-inj};
  \code{ToCentre}; \code{Defect} ($\mathsf{act}$, \code{pre-formula},
  Lemma~\ref{l:equiv} \code{inv-act}); \code{Diag} (Lemma~\ref{l:K}
  \code{kf-diag}, \code{mu}, \code{cell-mu}, Construction~\ref{c:R} \code{R},
  \code{wc}, \code{dprime}, Prop.~\ref{p:star} \code{cell-R},
  \code{ev-diag}, \code{inv-diag}); \code{PICircle} (Lemma~\ref{l:pi-circle});
  \code{repoint}; \code{DiagAt[B]}: \code{theta-inj}, Lemma~\ref{l:lift}
  (\code{Lift}, \code{lift-unique}), Lemma~\ref{l:phi} \code{phi-loop},
  Theorem~A (\code{diag-inj}, \code{undiag}, \code{diag-surj}).
\item \code{H4Bilinear}: \code{H2} (Lemma~\ref{l:H2}: \code{k1-inj},
  \code{k1-surj}, \code{k2-inj}, \code{k2-surj}); \code{EDef} ($\mathcal E$);
  \code{ENat} (Lemma~\ref{l:E-nat} pointwise: \code{nat-pt}, \code{hh},
  \code{e-nat}); \code{BilAt[B]}: \code{theta-inj} (Lemma~\ref{l:E-inj}),
  Lemma~\ref{l:j} \code{j-pt}, \code{key}, Theorem~B (\code{c2g-inj},
  \code{c2g-surj}).
\item \code{H4Proof}: \code{H4At[B]} (\code{inj}, \code{surj}, \code{equiv}),
  \code{h4}.
\end{itemize}
\end{document}
